\documentclass[11pt]{article}

% ---------- Packages ----------
\usepackage[utf8]{inputenc}
\usepackage[T1]{fontenc}
\IfFileExists{lmodern.sty}{\usepackage{lmodern}}{}
\usepackage[margin=1in]{geometry}
\usepackage{graphicx}
\usepackage{booktabs}
\usepackage{array}
\usepackage{tabularx}
\usepackage{ragged2e}
\usepackage{authblk}
\usepackage{abstract}
\usepackage{enumitem}
\usepackage[hidelinks]{hyperref}
\usepackage{caption}
\usepackage{xcolor}
\usepackage{textcomp}
\usepackage{pifont}

\captionsetup{font=small,labelfont=bf}
\renewcommand{\abstractnamefont}{\normalfont\bfseries}
\setlength{\parskip}{0.35em}

% ---------- Title ----------
\title{\textbf{AngioVision: An Open Source Multimodal Data Infrastructure for
Deep Angiography Research}}

% \author[1]{First A. Author\thanks{Corresponding author: author@institution.edu}}
% \author[1]{Second B. Author}
% \author[1,2]{Third C. Author}
% \affil[1]{Department of Radiology / Biomedical Informatics, Institution Name, City, Country}
% \affil[2]{Department of Computer Science, Institution Name, City, Country}

\date{}

\begin{document}
\maketitle

% =================== ABSTRACT ===================
\begin{abstract}
\noindent\textbf{Objectives:} To present AngioVision, an open-source,
web-accessible infrastructure that unifies DICOM metadata, free-text radiology
reports, and content-based image retrieval for natural-language and
image-similarity querying of angiography research archives.\\
\textbf{Materials and Methods:} The system comprises three components: a DICOM
metadata pipeline ingesting per-file tags into SQLite with SOPInstanceUID-based
deduplication; an image-embedding pipeline using the radiology-domain
self-supervised encoder RAD-DINO---with a selectable general-purpose Vision
Transformer baseline---indexing each encoder's vectors in a dedicated ChromaDB
collection; and a Flask web interface in which an agentic orchestrator
(smolagents) drives a locally hosted Qwen3 model through multi-step reasoning
over the relational metadata and report store and renders agent-generated charts
on request, with results streamed via server-sent events.\\
\textbf{Results:} The system architecture is described, including a
user-selectable image encoder (RAD-DINO or a general-purpose Vision Transformer)
with one vector collection per encoder. A reproducible evaluation framework
comprising 15 representative queries with blank scoring fields---including
image-retrieval items that support head-to-head encoder comparison---is provided
for institutional completion.\\
\textbf{Discussion:} The locally hosted design enables agentic,
retrieval-augmented querying without exposing protected health information,
and complements existing DICOM infrastructure.\\
\textbf{Conclusion:} AngioVision provides a reproducible, privacy-preserving
template for multimodal imaging-informatics research infrastructure.
\end{abstract}

\noindent\textbf{Keywords:} medical imaging informatics; DICOM; content-based
image retrieval; vector database; retrieval-augmented generation; LLM agents;
self-supervised learning; digital subtraction angiography

\vspace{0.5em}

% =================== MAIN TEXT ===================
\section{Background and Significance}

Digital subtraction angiography (DSA) remains the reference-standard modality
for vascular imaging, providing the high spatiotemporal resolution required for
diagnosis and image-guided endovascular
intervention~\cite{hill2007dsa,crummy2018history}. Deep-learning research on
angiographic data---vessel segmentation, artery--vein separation, hemodynamic
flow analysis~\cite{zhang2020nndsa}---increasingly requires large, well-curated
image archives linked to clinical context. In practice, such linkage is lacking:
research-scale archives consist of tens of thousands of DICOM objects whose
metadata are fragmented across files, whose pixel data are large and
multi-frame, and whose most informative content---the radiologist's
interpretation---resides in unstructured free text. The DICOM standard governs
storage and exchange~\cite{mildenberger2002dicom,bidgood1997dicom} but does not
provide flexible research-oriented query, and clinical PACS are designed for
routine workflow rather than exploratory cohort discovery.

Content-based image retrieval (CBIR) has been an active area of medical
imaging research for over two decades~\cite{litjens2017survey,muller2004cbir}.
Early CBIR systems relied on hand-crafted descriptors---color histograms,
texture features, shape moments---that captured low-level visual cues but
struggled with the semantic gap between pixel statistics and clinical
meaning~\cite{dubey2021cbir}. The advent of deep convolutional neural networks
transformed the field: features extracted from intermediate layers of
ImageNet-pretrained architectures (e.g., VGG, ResNet, DenseNet) proved
substantially more effective for medical image
retrieval~\cite{ke2021chextransfer,litjens2017survey}. However, natural-image
pretraining introduces a domain mismatch, as the learned features emphasize
textures and objects absent from radiographic
images~\cite{ke2021chextransfer}. Contrastive vision--language models such as
CLIP~\cite{radford2021clip} and its biomedical adaptations
BiomedCLIP~\cite{zhang2023biomedclip} and MedCLIP~\cite{wang2022medclip}
partially address this gap by aligning image and text representations within a
shared embedding space, enabling zero-shot classification and cross-modal
retrieval. Yet these models still require paired image--text data, which is
scarce in radiology owing to privacy constraints.

Self-supervised methods have emerged as a compelling alternative.
DINO~\cite{caron2021dino} and its successor DINOv2~\cite{oquab2024dinov2}
demonstrated that a self-distillation objective applied to Vision
Transformers~\cite{dosovitskiy2021vit} produces rich, general-purpose features
without any labels. RAD-DINO~\cite{perezgarcia2024raddino} extended this
paradigm to radiology, pretraining a ViT-B encoder on approximately 840{,}000
chest radiographs using DINOv2-style self-supervision. RAD-DINO outperformed
both image-only and image--text contrastive models across classification,
segmentation, and retrieval benchmarks, challenging the assumption that text
supervision is necessary for strong biomedical
encoders~\cite{perezgarcia2024raddino}. For efficient similarity search at
scale, deep hashing methods learn compact binary codes that preserve semantic
similarity~\cite{zheng2021ath}, while approximate nearest-neighbor (ANN)
indexes---in particular hierarchical navigable small-world (HNSW)
graphs~\cite{malkov2020hnsw}---enable sub-linear retrieval over dense
floating-point embeddings in vector databases.

Three advances thus converge to enable AngioVision. First, domain-specific
self-supervised encoders (RAD-DINO) provide powerful image representations
without paired text. Second, large language models (LLMs) translate natural
language into executable SQL~\cite{yu2018spider,pourreza2023dinsql} and ground
answers via retrieval-augmented generation
(RAG)~\cite{lewis2020rag,weinert2025rag}; agentic frameworks further allow
iterative, tool-using reasoning~\cite{yao2023react,wang2024codeact}. Third,
lightweight vector databases with ANN indexes make embedding-based retrieval
practical on commodity hardware. Although developed on an angiography archive,
the architecture is modality-agnostic: any radiologist, irrespective of
subspecialty or programming expertise, can deploy the system over their own
DICOM data.

\section{Objectives}

The objectives were to develop a privacy-preserving infrastructure that
(1)~ingests DICOM metadata and radiology reports into a queryable relational
store; (2)~computes radiology-domain image embeddings---with a selectable
general-purpose baseline---for content-based similarity retrieval; (3)~exposes
both through a web interface enabling
radiologists regardless of subspecialty or programming expertise to query their
archives via natural language and image similarity; and (4)~supports complex
multi-step questions through agentic reasoning. All model inference executes on
institutional hardware, ensuring no PHI leaves the
premises~\cite{thirunavukarasu2023llmmed,singhal2023clinical}.

\section{Materials and Methods}

\paragraph{System architecture.}
AngioVision pairs two managed data stores---SQLite for structured metadata and
free-text reports, and ChromaDB for image-embedding vectors---served by a single
Flask application (Figure~\ref{fig:arch}). Pixel data are not duplicated into a
separate store; individual frames are decoded on demand from the original DICOM
files, which are referenced by a \texttt{source\_path} pointer. Whichever image
encoder is selected is applied identically at ingestion and query time, so a
query embedding and the indexed vectors it is compared against always originate
from the same model and collection.

\paragraph{DICOM metadata pipeline.}
DICOM files are parsed via pydicom~\cite{mason2011pydicom} and stored as one row
per file in \texttt{dicom\_files}, spanning patient, study, series, and instance
identifiers; image geometry; radiation parameters (kVp, exposure time, tube
current, dose--area product); equipment attributes; and pipeline bookkeeping.
SOPInstanceUID-based deduplication avoids redundant re-parsing. A second table,
\texttt{radiology\_reports}, stores one row per accession number with the full
report text, linked without enforced foreign keys to accommodate incomplete data.

\paragraph{Image-embedding pipeline.}
Multi-frame pixel arrays are extracted via pydicom, corrected for MONOCHROME1
inversion, and normalized to uint8 RGB. By default each frame is embedded using
RAD-DINO (\texttt{microsoft/rad-dino}), a ViT-B encoder pretrained with
DINOv2-style self-supervision on approximately 840{,}000 chest
radiographs~\cite{perezgarcia2024raddino,oquab2024dinov2}. RAD-DINO was selected
for its radiology-domain specificity, replacing an earlier OpenCLIP
ViT-B/32 encoder~\cite{radford2021clip}. The encoder is pluggable: a
general-purpose ImageNet-pretrained ViT-B/16~\cite{dosovitskiy2021vit} is
provided as a selectable baseline, and further encoders can be added through a
single registry entry. For every encoder the CLS token (768-dimensional) is
$L_2$-normalized and stored in ChromaDB configured for cosine similarity over an
HNSW index~\cite{malkov2020hnsw}. Each encoder writes to its own ChromaDB
collection and records per-sequence ingestion state under its own model key, so
collections are built independently and compared head-to-head without
intermixing incomparable vector spaces. Image ingestion is single-process,
constrained by GPU memory and ChromaDB write semantics, and is resumable:
sequences already completed for a given encoder are skipped on re-run.

\paragraph{Web query interface.}
The Flask backend pairs with a dependency-free vanilla-JavaScript frontend
(HTML/CSS/JS, no build step; Figure~\ref{fig:ui}) and exposes three
capabilities: \textbf{(a)}~agentic natural-language querying via a locally hosted
Qwen3 model~\cite{yang2025qwen3} (selectable at query time, Qwen3-1.7B by
default) served through Ollama, whose answers may be accompanied by an
agent-generated chart---bar, line, pie/doughnut, or horizontal bar---drawn in
the browser by a dependency-free \texttt{<canvas>} renderer;
\textbf{(b)}~image-similarity retrieval, in which
a query image is embedded with the user-selected encoder, the corresponding
ChromaDB collection returns the top-$k$ similar sequences grouped by
SOPInstanceUID and accession number, and each hit is enriched with SQLite
metadata and a report excerpt; and \textbf{(c)}~on-demand frame extraction
handling MONOCHROME1 inversion, multi-frame indexing, and uint8 normalization.
Both query pathways stream progress---including any emitted chart
specification---to the browser via server-sent events (SSE).

\paragraph{Agentic query orchestration.}
Rather than a single-shot pipeline---one LLM call to generate SQL, up to two
automatic repairs, and a second call to synthesize an answer---query resolution
is delegated to an agentic orchestrator built on
smolagents~\cite{huggingface2025smolagents}. The model
follows a ReAct-style reason--act--observe loop~\cite{yao2023react}: at each
step it plans an action, invokes a tool, inspects the result, and iterates. The
agent is equipped with two tools, both issued through structured function calls
rather than free-form code. The first is a read-only SQL tool spanning the
structured metadata and the free-text reports---the latter searched with
\texttt{LIKE} predicates on the report column; with it the agent explores the
schema, samples distinct column values, recovers from SQL errors, and chains
multiple queries: for example, searching reports for a finding, extracting the
matching accession numbers, then aggregating the corresponding metadata. The
second is a chart-rendering tool: when a question calls for a visualization
(``plot'', ``distribution'', ``trend'', ``over time'', \ldots), the agent first
retrieves the underlying rows with the SQL tool, then invokes the chart tool
with a compact specification---chart type (bar, line, pie/doughnut, or
horizontal bar), title, and aligned category and value lists---which is streamed
to the browser and drawn client-side. Together these enable multi-hop
queries---and, where appropriate, their visual summaries---that a single-shot
pipeline cannot resolve. Content-based image retrieval is exposed as a
separate, deterministic endpoint rather than as an agent tool. All reasoning
executes locally on the selected Qwen3 model.

\paragraph{Implementation and availability.}
The stack comprises Python, Flask, a dependency-free vanilla-JavaScript frontend
(HTML/CSS/JS, no build step), pydicom, NumPy, Pillow, ChromaDB, the HuggingFace
\texttt{transformers}/PyTorch image encoders (RAD-DINO and a ViT-B/16 baseline),
Ollama serving the Qwen3 family, smolagents, and SQLite, deployed on a single
GPU-equipped Linux server. Command-line entry points (\texttt{run\_ingest.py} for
the ingestion pipelines and \texttt{run\_server.py} for the web interface) build
the data stores and launch the application, and a usage guide accompanies the
source. Source code is at \url{https://github.com/PLACEHOLDER/angiovision}.
Supplementary Table~S1 summarizes all components.

\section{Use Cases}

Three representative scenarios illustrate the system's workflow; all are
supported by the architecture in Figure~\ref{fig:arch}.

\paragraph{Use case 1: multi-step cohort assembly.}
An investigator enters: \emph{``Among DSA series whose report mentions an
aneurysm, which were acquired above 90~kVp, and what is the mean dose--area
product?''} The agentic orchestrator decomposes this into sequential tool calls:
report search, kVp filtering, and aggregation over the joined subset. If an
intermediate query fails, the agent revises its next action. The output includes
both the answer and the full query sequence for reproducibility.

\paragraph{Use case 2: content-based case retrieval.}
An investigator uploads a single DSA frame and requests visually similar cases
with reports. The frame is embedded with the selected encoder (RAD-DINO by
default), ChromaDB returns similar sequences grouped by accession number, and
each result is enriched with metadata
and a report excerpt---surfacing comparable cases that structured queries alone
would not identify, while all processing remains within the institutional
network.

\paragraph{Use case 3: visual summarization.}
An investigator asks: \emph{``Plot the number of studies acquired each year.''}
The agent runs a grouped SQL query, then calls the chart tool with the resulting
years and counts; a bar chart is streamed back and shown alongside the text
answer in the same result card. An aggregate query thus becomes an immediately
interpretable figure without exporting data to external plotting software.

\section{Results}

AngioVision constitutes data infrastructure rather than a predictive model;
accordingly, Table~\ref{tab:eval} defines 15 representative evaluation queries
spanning four categories: metadata (M), report text (R), complex multi-step (C),
and image similarity (I). Each query is paired with expected behavior and blank
scoring columns. For text-based queries, evaluators record query correctness,
execution success, number of agent steps, and answer quality (1--5 Likert). For
image queries, top-$k$ retrieval precision replaces query-level fields. Because
each encoder maintains an independent collection, the image-similarity items
(Q13--Q15) can be scored separately for RAD-DINO and the general-purpose
baseline, yielding a direct, site-specific retrieval comparison.
Aggregation yields per-category and composite scores; the blank template permits
site-specific reporting without disclosing protected data.

\section{Discussion}

AngioVision demonstrates that a locally hosted, multimodal infrastructure can
render a heterogeneous angiography archive queryable through both natural
language and image similarity. Grounding LLM output in executed queries and
retrieved metadata constitutes a form of
RAG~\cite{lewis2020rag} that mitigates hallucination in clinical LLM
applications~\cite{thirunavukarasu2023llmmed,weinert2025rag}. Because the agent
can additionally turn an aggregate query into a chart on request, common
descriptive analyses are completed entirely within the interface, without
exporting data to external plotting tools. The transition
from single-shot generation to an agentic
loop~\cite{huggingface2025smolagents,yao2023react} extends capabilities to
complex, multi-hop queries. The choice of RAD-DINO over general-purpose
encoders is supported by evidence that domain-specific self-supervised features
outperform both ImageNet-transferred and contrastive vision--language features
for radiology tasks~\cite{perezgarcia2024raddino,ke2021chextransfer}, while
the HNSW-backed vector database ensures sub-linear retrieval
latency~\cite{malkov2020hnsw}. Exposing the encoder as a user-selectable
component, with each model isolated in its own collection, turns this design
choice into a testable hypothesis: RAD-DINO and a general-purpose baseline can be
compared on identical archive data through the same interface.

\paragraph{Comparison with existing tools.}
Table~\ref{tab:comparison} positions AngioVision relative to established DICOM
infrastructure. Open-source platforms including
Orthanc~\cite{jodogne2018orthanc}, the OHIF Viewer~\cite{ziegler2020ohif}, and
DIANA~\cite{yi2021diana} provide robust storage, retrieval, and viewing;
commercial suites such as mPower add enterprise-scale analytics. None offers
natural-language metadata querying, agentic multi-step reasoning, or
embedding-based image retrieval. AngioVision is designed as a complementary
query layer over data such archives already manage well. Its deliberate omission
of 3D/MPR rendering constitutes the principal trade-off; pairing with OHIF
addresses this limitation.

\paragraph{Generalizability.}
The natural-language interface eliminates the need to write SQL or Python; the
agentic orchestrator decomposes complex questions on behalf of the user; and the
image-similarity endpoint requires only a single uploaded frame. Neither the
ingestion pipelines, the embedding encoders, nor the query interface contain
angiography-specific logic. Replacing the data source requires only re-running
ingestion on a new DICOM root. RAD-DINO, trained on a broad radiographic
corpus, generalizes across modalities~\cite{perezgarcia2024raddino}; and because
the image encoder is user-selectable with each model indexed in its own
collection, alternative encoders can be benchmarked head-to-head on the same
archive without re-architecting the system.

\paragraph{Limitations.}
Image ingestion is constrained to single-process execution. Referential
integrity between studies and reports is not enforced. The interface inherits
text-to-SQL error modes documented for
LLMs~\cite{pourreza2023dinsql,yu2018spider}. RAD-DINO, trained predominantly on
chest radiographs, may require domain adaptation for angiographic sequences;
the selectable general-purpose ViT-B/16 baseline shares ImageNet's domain gap,
and a quantitative head-to-head retrieval comparison between the two on
angiographic data remains future work. Formal site-specific evaluation using
Table~\ref{tab:eval} is an essential next step.

\section{Conclusion}

AngioVision is an open-source, reproducible infrastructure template for
multimodal imaging informatics that integrates DICOM metadata, radiology
reports, and domain-specific image embeddings under a unified agentic
natural-language and image-similarity query interface. All inference executes
locally, preserving patient privacy. Although developed on an angiography
archive, the architecture is modality-agnostic: any radiologist, regardless of
subspecialty or programming expertise, can deploy and query the system over
their own DICOM data without modifying the core pipeline.

% =================== FIGURES ===================
\begin{figure}[ht]
  \centering
  \includegraphics[width=0.95\textwidth]{figs/fig1_architecture.png}
  \caption{\textbf{System architecture.} Two managed stores---SQLite
  (\texttt{dicom\_files}, \texttt{radiology\_reports}) and ChromaDB
  ($L_2$-normalized 768-dimensional CLS embeddings, cosine HNSW index)---are
  populated by two ingestion pipelines and queried through a Flask web
  interface; pixel data are decoded on demand from the original DICOM files
  (\texttt{source\_path}). An agentic orchestrator (smolagents) drives a locally
  hosted Qwen3 model (Ollama; selectable, Qwen3-1.7B by default) through
  iterative reason--act--observe loops over a read-only SQL tool and a
  chart-rendering tool. The
  selected image encoder---RAD-DINO by default, or a general-purpose ViT-B/16
  baseline---is used identically at ingestion and query time, with one ChromaDB
  collection per encoder. All inference executes on a single GPU-equipped Linux
  server; no PHI leaves the premises.}
  \label{fig:arch}
\end{figure}

\begin{figure}[ht]
  \centering
  \includegraphics[width=0.92\textwidth]{figs/fig2_interface.png}
  \caption{\textbf{Web query interface.} The vanilla-JavaScript frontend
  showing a natural-language query with streamed answer and generated SQL
  (left), and image-similarity results (right): visually similar DICOM sequences
  grouped by SOPInstanceUID and accession number, enriched with metadata and
  radiology-report excerpts. A sidebar control selects the active image encoder
  and reports the size of its vector collection. The interface supports
  progressive thumbnail loading, a full-resolution lightbox, and accession-level
  deduplication.}
  \label{fig:ui}
\end{figure}

% =================== TABLES ===================
% --- Supplementary (unnumbered) component table ---
\begin{table}[hb]
\centering
\caption*{\textbf{Supplementary Table S1.} AngioVision components, functions,
and key dependencies.}
\label{tab:components}
\small
\renewcommand{\arraystretch}{1.2}
\begin{tabularx}{\textwidth}{@{}l X X@{}}
\toprule
\textbf{Component} & \textbf{Function} & \textbf{Dependencies} \\
\midrule
DICOM metadata pipeline &
One row per \texttt{.dcm} in \texttt{dicom\_files}; SOPInstanceUID-based
deduplication; reports keyed by accession number &
pydicom, SQLite \\[0.4em]
Image-embedding pipeline &
Multi-frame extraction, uint8 RGB normalization, CLS embedding
($L_2$-normalized, cosine); selectable encoder (RAD-DINO or ViT-B/16), one
collection per encoder; resumable, single-process &
pydicom, NumPy, Pillow, RAD-DINO / ViT-B, ChromaDB \\[0.4em]
Agentic query interface &
Multi-step reason--act--observe loop over a read-only SQL tool (metadata +
free-text reports) and a chart-rendering tool; SSE streaming &
Flask, smolagents, Ollama (Qwen3 family), SQLite \\[0.4em]
Image-similarity interface &
Embed query image via the selected encoder; per-encoder ChromaDB top-$k$; group
by SOPInstanceUID/accession; SQLite enrichment &
Flask, RAD-DINO / ViT-B, ChromaDB, SQLite \\[0.4em]
Frame-extraction endpoints &
On-demand thumbnails and full-resolution frames; MONOCHROME1 inversion;
multi-frame handling &
Flask, pydicom, NumPy, Pillow \\[0.4em]
Frontend &
Vanilla JS (HTML/CSS/JS, no build step); progressive loading; lightbox;
self-contained \texttt{<canvas>} chart rendering; accession-level deduplication &
JavaScript \\
\bottomrule
\end{tabularx}
\end{table}

% --- Table 1: Evaluation framework ---
\begin{table}[ht]
\centering
\caption{\textbf{Evaluation framework.} Representative queries with blank
scoring fields for institutional completion. Q-type:
M\,=\,metadata, R\,=\,report text, C\,=\,complex/multi-step (agentic),
I\,=\,image similarity. Steps\,=\,number of agent tool calls.
For image queries, retrieval relevance (top-$k$ precision) replaces query-level
fields.}
\label{tab:eval}
\scriptsize
\renewcommand{\arraystretch}{1.25}
\begin{tabularx}{\textwidth}{@{}c l c X c c c c@{}}
\toprule
\textbf{ID} & \textbf{Query} & \textbf{Type} &
\textbf{Expected behavior} & \textbf{Correct (Y/N)} &
\textbf{Exec.} & \textbf{Steps} & \textbf{Quality (1--5)} \\
\midrule
Q1 & Total distinct studies in the archive & M &
\texttt{COUNT(DISTINCT study\_instance\_uid)} & & & & \\
Q2 & Series acquired with kVp $>$ 90 & M &
Filter on \texttt{kvp > 90}; return series & & & & \\
Q3 & Mean dose--area product per accession & M &
Aggregate \texttt{dose\_product} by \texttt{accession\_number} & & & & \\
Q4 & Frame count per multi-frame instance & M &
Return \texttt{frame\_count} per \texttt{sop\_instance\_uid} & & & & \\
Q5 & Equipment manufacturers and frequencies & M &
\texttt{GROUP BY manufacturer} with counts & & & & \\
Q6 & Studies with report mentioning ``aneurysm'' & R &
Join; \texttt{radrpt LIKE '\%aneurysm\%'} & & & & \\
Q7 & Reports describing carotid stenosis & R &
Text match in \texttt{radrpt}; return accessions & & & & \\
Q8 & Summarize findings for accession X & R &
Retrieve report; synthesize summary & & & & \\
Q9 & Count of accessions with normal angiogram & R &
Text match; \texttt{COUNT} & & & & \\
Q10 & Mean dose--area product of $>$90\,kVp series among
aneurysm-positive reports & C &
Report search $\to$ kVp filter $\to$ aggregate (multi-step) & & & & \\
Q11 & Per-manufacturer count of stenosis-positive studies & C &
Join reports to equipment; group; count & & & & \\
Q12 & Year with most high-dose studies; summarize those reports & C &
Aggregate by year $\to$ identify peak $\to$ summarize & & & & \\
\midrule
\multicolumn{4}{@{}l}{\textit{Image-similarity queries}} &
\multicolumn{4}{c}{\textbf{Retrieval relevance / Quality}} \\
Q13 & (Upload DSA frame) Find similar sequences & I &
Embed; top-$k$; group by accession & & & & \\
Q14 & Cases resembling this aneurysm projection & I &
Top-$k$ with report enrichment & & & & \\
Q15 & Prior studies for the same anatomy & I &
Retrieval + enrichment; accession dedup & & & & \\
\bottomrule
\end{tabularx}
\end{table}

% --- Table 2: Feature comparison ---
\begin{table}[ht]
\centering
\caption{\textbf{Feature comparison with existing DICOM tools.}
\ding{51}\,=\,supported; \ding{55}\,=\,not supported; P\,=\,partial.
mPower is a commercial platform (Bayer/Nuance);
Orthanc~\cite{jodogne2018orthanc}, dcm4chee\,+\,OHIF~\cite{ziegler2020ohif},
and DIANA~\cite{yi2021diana} are open-source. Entries reflect each system's
primary design; some capabilities may be achievable via plugins.}
\label{tab:comparison}
\footnotesize
\renewcommand{\arraystretch}{1.3}
\setlength{\tabcolsep}{4pt}
\begin{tabularx}{\textwidth}{@{}l X c c c c c@{}}
\toprule
\textbf{Category} & \textbf{Capability} &
\textbf{\shortstack{Angio-\\Vision}} & \textbf{mPower} &
\textbf{Orthanc} & \textbf{\shortstack{dcm4chee\\+\,OHIF}} &
\textbf{DIANA} \\
\midrule
Access \& cost
 & Open source & \ding{51} & \ding{55}$^{\mathrm{a}}$ & \ding{51} & \ding{51} & \ding{51} \\
 & On-premise data residency & \ding{51} & \ding{55} & \ding{51} & \ding{51} & \ding{51} \\
\addlinespace[2pt]
Data \& ingestion
 & High-throughput ingestion$^{\mathrm{b}}$ & \ding{51} & \ding{51} & P & \ding{51} & P \\
 & Crash-safe ingestion / audit & \ding{51} & \ding{55} & \ding{51} & \ding{51} & \ding{51} \\
 & Missing AccessionNumber handling & \ding{51} & \ding{55} & \ding{55} & \ding{55} & \ding{55} \\
\addlinespace[2pt]
Query interface
 & NL metadata queries & \ding{51} & P$^{\mathrm{c}}$ & \ding{55} & \ding{55} & \ding{55} \\
 & DSA-specific metadata queries & \ding{51} & \ding{55} & P & P & P \\
 & Agentic multi-step reasoning & \ding{51} & \ding{55} & \ding{55} & \ding{55} & \ding{55} \\
 & Content-based image similarity & \ding{51} & \ding{55} & \ding{55} & \ding{55} & \ding{55} \\
\addlinespace[2pt]
Viewing
 & Image viewing / 3D / MPR & \ding{55}$^{\mathrm{d}}$ & \ding{55} & P & \ding{51} & P \\
\bottomrule
\end{tabularx}

\vspace{2pt}
{\footnotesize
$^{\mathrm{a}}$Commercial SaaS.\quad
$^{\mathrm{b}}$Reference: ${\sim}$816{,}000 files in ${\sim}$35\,min.\quad
$^{\mathrm{c}}$Report-text NLP only.\quad
$^{\mathrm{d}}$Metadata-and-retrieval layer; designed to pair with OHIF or
similar viewers.}
\end{table}

% =================== BIBLIOGRAPHY ===================
\bibliographystyle{unsrt}
\bibliography{references}

\end{document}